\chapter{Implementation Details}
\label{ch:implementation}

\section{Development Environment}
To set up the project and begin writing and running code, I configured a local development environment. The tools and software versions I used during my practical training are listed in Table~\ref{tab:development-environment}.

\begin{table}[htbp]
  \centering
  \caption{My practical training development environment}
  \label{tab:development-environment}
  \begin{tabularx}{\textwidth}{p{0.30\textwidth} X}
    \toprule
    \textbf{Component} & \textbf{Technology / Tool Used}\\
    \midrule
    Operating System & macOS (Darwin) and Windows 11 \\
    Programming Languages & Java 21, Dart (Flutter 3), Python 3.10, TypeScript (React 19)\\
    Frameworks \& Libraries & Spring Boot 3, Spring Security, FastAPI, HuggingFace Transformers, PyTorch\\
    Databases & PostgreSQL 15, Flutter SharedPreferences\\
    Containerization & Docker and Docker Compose v2\\
    IDE \& Tooling & VS Code, IntelliJ IDEA, Maven, Git\\
    \bottomrule
  \end{tabularx}
\end{table}

\section{Monorepo Project Structure}
I worked inside a single, structured monorepo containing all the services. To manage the code, I focused on these four main folders:
\begin{enumerate}
  \item \textbf{\texttt{frontend/}:} This contains the canonical Flutter mobile app (named `secure_signal`). It has its own \texttt{lib/main.dart} file, customized views for login/scans, and the background SMS listener.
  \item \textbf{\texttt{backend/}:} The Spring Boot Java API, which contains Maven configuration (\texttt{pom.xml}), Flyway database migrations (\texttt{db/migration}), and java controllers.
  \item \textbf{\texttt{web-client/}:} The administrator dashboard built with React and Vite.
  \item \texttt{main.py} and \textbf{\texttt{final\_model/}:} The Python script that runs FastAPI and the folder containing the saved BERT model weights and config files.
\end{enumerate}

Note: The root \texttt{lib/} directory is a legacy or older version of the mobile app, which we retired to avoid duplication, making \texttt{frontend/} our main focus.

\section{Data Preprocessing and Heuristics Implementation}
The local rule-based engine in Flutter uses regular expressions and case-insensitive word matching. For example, if a message matches words like "AirtelMoney", "M-Pesa", "win prize", "transfer to", or Swahili phrases such as "utatuma" (you will send) or "hakikisha jina", the app immediately identifies the threat level.

On the machine learning service, the input string must be checked to make sure it is not empty. If the message is valid, it is prepared for BERT using the HuggingFace tokenizer. The implementation is shown in Listing~\ref{lst:preprocessing-excerpt}.

\begin{lstlisting}[language=Python,caption={My FastAPI text prediction endpoint implementation},label={lst:preprocessing-excerpt}]
@app.post("/predict")
def predict(request: PredictRequest):
    # Make sure message is not empty
    if not request.message.strip():
        raise HTTPException(status_code=400, detail="message is required")
        
    # Tokenize the text to fit the BERT model's 128 limit
    encoded = tokenizer(request.message, truncation=True,
                        max_length=128, return_tensors="pt")
                        
    # Run the model to get raw prediction outputs
    output = model(**encoded)
    
    # Format and return the label and confidence percentage
    return format_label_and_confidence(output, request.message)
\end{lstlisting}

This API code is simple, reliable, and runs on port 8000, allowing our Spring Boot service to connect to it over the local network.

\section{Model Loading and Prediction}
When the FastAPI application starts, the BERT classifier is loaded from the directory \texttt{./final\_model} using the \texttt{from\_pretrained()} method from the HuggingFace library. The model weights are loaded into CPU memory (or CUDA GPU if available). Since the model is loaded once at startup, predictions are extremely fast, usually taking only a few milliseconds.

\section{Backend and Database Integration}
The Spring Boot backend uses Flyway to automatically create and update database tables on startup. I analyzed the database migration scripts:
\begin{itemize}
  \item \texttt{V1__init.sql}: Creates tables for users, roles, and admin credentials.
  \item \texttt{V3__sms_scans.sql}: Creates the \texttt{sms_scans} table to log malicious messages.
  \item \texttt{V4__make_user_id_and_verdict_nullable.sql}: Simplifies integration by allowing scan submissions from anonymous or unregistered users.
\end{itemize}

The Spring Boot backend acts as a bridge. When the mobile app requests a scan, the backend routes the message to the Python FastAPI, receives the result, and if the message is a scam, records the sender, body, confidence, and timestamp into the PostgreSQL database.

\section{Deployment with Docker Compose}
To make testing easier, I packaged the entire stack using Docker Compose. The \texttt{docker-compose.yml} file launches:
\begin{enumerate}
  \item A PostgreSQL 15 database on port 5432, with a persistent folder to make sure our data is not lost when we stop the container.
  \item The Python ML service on port 8000, which loads the saved BERT model.
  \item The Spring Boot backend on port 8080, which connects to both PostgreSQL and the ML service using Docker's internal networking (\texttt{http://ml-service:8000}).
\end{enumerate}

Running \texttt{docker-compose up} automatically builds, links, and starts all these services, which saved me a lot of time during my practical training.

\section{Implementation Challenges}
During the integration phase of my practical training, I encountered several real-world software engineering challenges. I have summarized these challenges, their effects, and how I resolved them in Table~\ref{tab:implementation-challenges}.

\begin{table}[htbp]
  \centering
  \caption{Implementation challenges and solutions}
  \label{tab:implementation-challenges}
  \begin{tabularx}{\textwidth}{p{0.28\textwidth} p{0.32\textwidth} X}
    \toprule
    \textbf{Challenge} & \textbf{Effect on Project} & \textbf{My Solution / Decision}\\
    \midrule
    Two Flutter folders & Created confusion about which codebase was the active one. & I chose the \texttt{frontend/} directory as the main code and retired the root folder.\\
    Inconsistent login endpoints & Web and mobile apps expected different URLs for login. & I updated the Dart and React code to use the new resource paths (\texttt{/api/auth/...}).\\
    CORS and Network Blockers & The browser blocked React requests due to security rules. & I configured a CORS filter in Spring Boot to allow connections from the React port.\\
    Different database models & Web dashboard was using mock arrays while backend had real API logic. & I documented the database schemas to help future developers connect the live API.\\
    \bottomrule
  \end{tabularx}
\end{table}

\section{Chapter Summary}
In this chapter, I documented the development environment and structural layout of the monorepo. I explained how the heuristic and BERT models are implemented, how Spring Boot communicates with PostgreSQL and FastAPI, and how Docker Compose packages the services. I also listed the integration challenges I faced and how I solved them. In the next chapter, I will cover the testing and results of the system.
